% RESEARCH_DRAFT_SHA256: 4e08772dcb9c4443e752651fad46fd57d1782281f85c4a9bb2ba31013cf92c3a
% Layout and source-reference wrappers expanded for the web renderer.
\documentclass[11pt]{article}
\usepackage[margin=25mm]{geometry}
\usepackage{fontspec,unicode-math}
\setmainfont{Latin Modern Roman}
\setmathfont{Latin Modern Math}
\usepackage{amsmath,amssymb,mathtools,xurl,hyperref}
\hypersetup{colorlinks=true,urlcolor=blue,linkcolor=blue}
\setlength{\parindent}{0pt}
\setlength{\parskip}{5pt}
\setlength{\emergencystretch}{3em}
\newcommand{\sref}[2]{\hyperlink{source:#1:#2}{[S#2]}}
\newcommand{\cataloguescope}[1]{\paragraph{Source question.}#1}
\newcommand{\floor}[1]{\left\lfloor#1\right\rfloor}
\begin{document}

\section*{Pooling hull fragment: a recovered model and an explicit failure of separate convexification}
\noindent catalogue entry; catalogue entry. Research attempt dated 11 October 2026.
\paragraph{Result in scope.}This manuscript contains a source-specific attempt and complete proofs of the stated auxiliary calculations. It makes no new solution or novelty claim for the entire target.
\section{Original problem and preserved definition}
The following definition and source list are reproduced from the repository. Any earlier inspection dates and reported literature claims in this preserved text describe that source record, not additional verification by this attempt.
% BEGIN PRESERVED DEFINITION
\subsection{Definitions and mathematical statement}

\paragraph{Shared notation.}Standard number systems are $\mathbb N,\mathbb Z,\mathbb Q,\mathbb R,\mathbb C$. Specialized notation and index ranges follow the cited source and the definitions below.


The convex hull of a set consists of its finite convex combinations. A pooling model mixes quantities while imposing consistency between flow amounts and their blended qualities; these conditions can be nonconvex.
\paragraph{Record scope and qualifications.}
The catalogue title is retained as metadata. Named certificates, contracts, bounds or frameworks in that title are not asserted as results of the original source extract.
The supplied registry flags this extraction as unrecoverable. It is displayed as a literal source fragment; the definitions provide context rather than a replacement formal statement.

The underlying nonconvex set and the proposed hull inequalities are absent.

Source-record limitation: Source flags the statement as unrecoverable: The official report only says that a structured nonconvex pooling relaxation has a conjectured full convex-hull description in the remaining cases; it does not state the subset, cases, or conjectured inequalities. No separate clean statement is promoted, and the original record remains visible for provenance.
The source text below is rendered literally where necessary; damaged formulas and composite extracts are not silently repaired into a new mathematical claim.
The supplied catalogue marks this record as unrecoverable. Its extract and terminology are documented, but a complete original problem has not been reconstructed.
\paragraph{Statement recorded in the supplied source.}
For the remaining considered cases of the structured nonconvex subset of the pooling problem, is the conjectured full description of its convex hull correct? \href{https://doi.org/10.4171/owr/2015/46}{[S2]}
\paragraph{Source status.}
The supplied catalogue labels this source record open; that metadata was checked against the supplied snapshot on 2026-10-11. The label is preserved for provenance and is not an independent verification of present open status or mathematical correctness. The catalogue literature assessment is dated 2026-08-21.
\subsection{Short English statement}
The fragment asks whether a proposed pooling convex-hull description is complete.
\subsection{Sources}
\begin{itemize}
\item[S1] Ulam AI and the UnsolvedMath Contributors, UnsolvedMath: A Curated Collection of Open Mathematics Problems, catalogue entry. Catalogue curation and metadata: CC BY 4.0. The source extract is retained; definitions, English summaries and qualifications are editorial additions.
\url{https://huggingface.co/datasets/ulamai/UnsolvedMath}.
\item[S2] Mixed-integer Nonlinear Optimization: A Hatchery for Modern Mathematics (2016), Oberwolfach Reports, DOI 10.4171/owr/2015/46, pp. 2701--2766
\url{https://doi.org/10.4171/owr/2015/46}.
\item Source-specific Oberwolfach Reports licences, including CC BY-SA 4.0 where applicable, are retained; see THIRD\_PARTY\_NOTICES.md.
\end{itemize}
% END PRESERVED DEFINITION
\subsection{Verbatim supplied catalogue question}
For the remaining considered cases of the structured nonconvex subset of the pooling problem, is the conjectured full description of its convex hull correct?

\section{Source limitation and later primary model}
The OWR passage, printed pages 2715--2716, does not list the parameters,
remaining cases, or proposed inequalities. The source record therefore stays
qualified as unrecoverable. A primary-source search found the later paper by
Luedtke, D'Ambrosio, Linderoth and Schweiger, arXiv:1803.02955, published in
SIAM Journal on Optimization in 2020. Its introduction and Sections 2--3
were inspected. It describes a five-variable aggregate model and reports
full hull descriptions in three parameter cases. The complete proofs of those
three cases and their exact correspondence to the unnamed 2015 remaining
cases have not been audited. No full solution is inferred from the abstract.

We can nevertheless test a source-specific mathematical route on an explicitly
specified specialization of its displayed model. In the later paper take
capacity $C=1$, pool-quality bounds $0\le t\le1$, and bypass-quality bounds
$-1\le y/z\le0$, interpreted linearly at $z=0$. The set is
\[
T=\{(x,u,y,z,t):0\le x\le1,\ 0\le t\le1,\ z\ge0,\ x+z\le1,
-z\le y\le0,\ y+u\le0,\ u=xt\}.
\]
This specialization is an auxiliary test, not a silently substituted statement
of the missing original conjecture.

\section{The isolated bilinear hull and the failed intersection step}
For the graph $G=\{(x,t,u):0\le x,t\le1,u=xt\}$, its convex hull is exactly
\[
0\le x,t\le1,\qquad
u\ge0,\quad u\ge x+t-1,\quad u\le x,\quad u\le t.
\]
Validity on the graph follows from the bounds. Conversely any point satisfying
these inequalities is the convex combination of the four graph corner points
$(0,0,0),(1,0,0),(0,1,0),(1,1,1)$ with respective weights
\[
1-x-t+u,\quad x-u,\quad t-u,\quad u.
\]
They are nonnegative, sum to one, and reproduce all three coordinates.
This proves the isolated hull with no external envelope theorem.

An attempted hull description for $T$ replaces only $u=xt$ by those four
inequalities and retains the other linear constraints. It fails at
\[
(x,u,y,z,t)=(1/2,1/2,-1/2,1/2,1/2).
\]
All linear side constraints and all four bilinear-hull inequalities hold.
We next derive a convex inequality valid on $T$ that excludes this point,
so it is not even in $\operatorname{conv}T$.

\section{A complete nonlinear certificate}
For a point of $T$, $u=xt\ge0$ and
$u\le-y\le z\le1-x$, so $u+x\le1$.
Multiplication by $u\ge0$ gives
\[
u(u+x)\le u=xt,\qquad\text{equivalently}\qquad
u^2\le x(t-u).
\]
Here $x\ge0$ and $t-u=t(1-x)\ge0$. The set
$\{(x,t,u):x\ge0,t-u\ge0,u^2\le x(t-u)\}$ is convex.
Indeed for any finite convex combination of feasible triples, write
$a_i=x_i,b_i=t_i-u_i$ and weights $\theta_i$. Then
\[
\left|\sum_i\theta_i u_i\right|
\le\sum_i\theta_i\sqrt{a_i b_i}
\le\sqrt{\left(\sum_i\theta_i a_i\right)
\left(\sum_i\theta_i b_i\right)}
\]
by Cauchy--Schwarz, establishing the quadratic inequality and the required
nonnegative sides for the combination. Thus it is valid on the whole hull.
At the proposed point, $u^2=1/4$ but $x(t-u)=0$, so the inequality fails.
This is the corresponding simple specialization of the later paper's
nonlinear cut, rederived here with a full direct proof.

The mathematical error isolated by the example is
$\operatorname{conv}(G\cap L)=\operatorname{conv}(G)\cap L$ for an arbitrary
linear side-constraint set $L$: the four-corner representation may use
components that violate $L$, even when their average meets $L$.
Exact convexification of each product separately therefore does not prove a
full pooling hull. The one additional cut is shown valid and necessary for
this example; its sufficiency for this or every other parameter case is not
claimed or proved.

\section{Adversarial audit and remaining scope}
The audit checked the bypass sign bounds at zero flow, convexity rather than
mere validity on the nonconvex set, and every coordinate of the fractional
witness. The new primary lead improves identification of the relevant model,
but does not restore the omitted 2015 quantifiers automatically. The result
is a known-cut rederivation and a failed-route certificate, not a new hull
theorem or a resolution of the unspecified original cases. No numerical
computation, theorem prover or independent human proof checking was used.

\section{Process and contributions}
Mathematics and model-assisted adversarial self-audit: OpenAI gpt-6-astra, ultra reasoning, as dispatched for this research batch; exact serving revision not exposed. Codex (GPT-6; exact serving revision not exposed) materially assisted with orchestration and format review. No human mathematical authorship, mathematical edits, or independent human proof checking is asserted.
The coordinating Codex (GPT-6; exact serving revision not exposed) also provided a separate model-assisted written-proof review of this manuscript. This is not independent human verification.
Contribution prompt A was applied to the full definition and supplied source row, then prompt B was applied adversarially to the derived results. The item-specific audit above records the scope checks and rejected extrapolations. Local assembly and structural checks do not establish mathematical correctness. The source catalogue status is preserved in TOP\_PROBLEM and is separate from this attempt status.
\paragraph{Integration review.} Codex (GPT-6; exact serving revision not exposed) reviewed the written arguments and expanded layout and source-reference wrappers. Repository metadata and the production web renderer were checked using Python 3.13 and JavaScript with MathJax 3.2.2. This model-assisted review is not independent human proof verification; no PDF compilation or formal proof checking is asserted.

\section{Completion Estimate}
\noindent\textbf{COMPLETION: 0\%}
The proved material is a rederivation of known or auxiliary facts rather than a substantiated new advance on the original target. The unnamed original cases and full proposed hull remain unverified; a later model was recovered and one valid nonlinear cut plus a counterexample to separate convexification was proved.
\section{References}
The complete original Sources subsection above is retained, including catalogue attribution and licence notices. The additional list records primary sources actually consulted for this attempt on 11 October 2026; uncited source links reproduced above were not all rechecked.
\begin{itemize}
\item Oberwolfach Report 46/2015, open-problem session pp. 2707--2708 and the pooling discussion pp. 2715--2716; exact passages inspected in the primary report. \url{https://doi.org/10.4171/owr/2015/46}.
\item Luedtke, D’Ambrosio, Linderoth and Schweiger, Strong Convex Nonlinear Relaxations of the Pooling Problem; introduction and Sections 2--3, including the model and first nonlinear inequality, inspected; full three-case hull proofs not checked. \url{https://arxiv.org/html/1803.02955}.
\item Publisher abstract and publication metadata for the 2020 SIAM Journal on Optimization paper inspected. \url{https://epubs.siam.org/doi/10.1137/18M1174374}.
\end{itemize}
\section{Final status}
\textbf{Unresolved in this attempt.} The unnamed original cases and full proposed hull remain unverified; a later model was recovered and one valid nonlinear cut plus a counterexample to separate convexification was proved.
\end{document}
